\documentclass[11pt,letterpaper]{article}
\usepackage[margin=1.02in,top=.91in,bottom=.91in,headheight=14pt]{geometry}
\usepackage[T1]{fontenc}\usepackage{lmodern}
\usepackage{amsmath,amssymb,amsthm,booktabs,array,xcolor,fancyhdr,hyperref}
\definecolor{ink}{HTML}{153950}\definecolor{accent}{HTML}{116C70}
\hypersetup{colorlinks=true,urlcolor=accent,linkcolor=accent,citecolor=accent,pdftitle={When does the spanning-tree formula start?}}
\pagestyle{fancy}\fancyhf{}\fancyhead[L]{\small\color{ink}Integral voltage towers}\fancyhead[R]{\small\color{ink}Torsion specializations}\fancyfoot[C]{\small\thepage}
\renewcommand{\headrulewidth}{.3pt}
\newtheorem{theorem}{Theorem}\newtheorem{lemma}[theorem]{Lemma}\newtheorem{corollary}[theorem]{Corollary}
\newtheorem{proposition}[theorem]{Proposition}
\newcommand{\Z}{\mathbb Z}\newcommand{\Q}{\mathbb Q}\newcommand{\F}{\mathbb F}
\newcommand{\leanfile}[2]{\href{https://github.com/lsernau/grh-lean/blob/72fead102f663d5c7a407964a13dc0cf9d977207/grh/src/#1}{#2}}
\setlength{\parskip}{3pt}\setlength{\parindent}{1.15em}
\setlength{\emergencystretch}{2em}
\setcounter{tocdepth}{1}\raggedbottom
\begin{document}
\thispagestyle{plain}
\noindent{\sffamily\small\color{accent} RESEARCH NOTE \hfill 26 NOVEMBER 2026}\par
\vspace{.15in}
\noindent{\LARGE\bfseries\color{ink} When does the spanning-tree formula start?}\par
\vspace{3pt}
\noindent{\large Effective Iwasawa polynomials from cyclotomic valuations}\par
\vspace{.16in}
\noindent\textbf{Abstract.}
A tower of connected graphs with deck groups $(\Z/p^n\Z)^d$ has a spanning-tree valuation eventually given by a polynomial in $p^n$ and $n$. We give an algorithm that takes integral voltage vectors, returns the polynomial, and certifies the last exceptional level. This answers the first research question in the ProveIt graph-tower report~\cite{incoming}. The construction combines the cyclotomic valuation calculation used in the Lean proof of Dirichlet GRH with a finite-state computation of binomial coefficients. It applies in every rank. Five loops suffice at $p=3$ to prescribe any positive starting level. A second family shows that even the leading coefficient cannot be computed uniformly from digit-oracle voltages. An application to Grover walks gives an all-level determinant formula on the ternary square torus.

\section{The problem and its scope}
A voltage is a vector attached to an oriented edge of a base graph. It gives the displacement between the copies of its endpoints in a covering graph. For a bouquet (one vertex and $m$ loops), give the loops steps $s_1,\ldots,s_m\in\Z^d$. At level $n$, the vertices are $G_n=(\Z/p^n\Z)^d$. For every $x$ and every listed step put one undirected edge between $x$ and $x+s_i$. Repetitions count. Assume the steps span $\F_p^d$ after reduction. They then generate every $G_n$. Write $\kappa_n$ for the number of spanning trees and $T_n=v_p(\kappa_n)$, normalized by $v_p(p)=1$. We include $n=0$, so $T_0=0$ for a bouquet.

\begin{theorem}[Effective polynomial and first valid level]\label{thm:main}
Given a prime $p$ and a finite list of integer steps as above, one can compute the unique rational polynomial $P(X,Y)$ of total degree at most $d$, of degree at most one in $Y$, and the least $N\ge0$ such that $T_n=P(p^n,n)$ for every $n\ge N$. The algorithm also returns the exact rational generating function of all $T_n$ and the finite polynomial of exceptions. If levels start at one, the least start is $\max(1,N)$.
\end{theorem}
DuBose and Vallières prove existence in their Theorem 6.2 and compute finite examples in their \S7~\cite{DV}. A prior ProveIt report proves their five candidate formulas and asks for a terminating procedure for arbitrary rank-two integral steps~\cite{incoming}. Theorem~\ref{thm:main} answers that algorithmic question. The same construction handles connected multi-vertex base graphs and rational voltages with denominator prime to $p$.

\section{What comes from the Lean proof}
The repository at commit \texttt{72fead102f66} states RH for the Riemann zeta function and GRH for Dirichlet $L$-functions. Its public statements use Mathlib's zeta function and Dirichlet characters, including imprimitive characters. The repository reports a Lean 4.35.0 build and proof-export checks, with \texttt{propext}, \texttt{Classical.choice}, and \texttt{Quot.sound} as the public theorems' axioms. We inspected the declarations and proof text. The upstream build and export checks were not repeated here. \cite{repo}

In the main argument, short integral forms and a spanning estimate give an upper bound for a corrected covolume. A lower bound from twisted determinants controls the same quantity. Subtracting common terms leaves an inequality for a weighted prime sum. A one-sided estimate and a choice of quadratic twist then exclude zeros. A transfer to a new family of $L$-functions would need a corresponding prime statistic, an estimate of the requisite strength, and control of cancellations among residues. The source establishes these inputs for its Dirichlet-character setup. See \leanfile{New/Reduction/MainInequality.lean}{the main inequality}, \leanfile{New/Reduction/OneSidedBound.lean}{the one-sided estimate}, and \leanfile{New/Reduction/Zeros/Criterion.lean}{the zero criterion}.

Another part of the proof works directly with polynomials at roots of unity. In \leanfile{New/Congruences/Selmer/RootsOfUnity.lean}{the Selmer specialization file}, the lemma \texttt{ellVal\_sub\_one\_eq\_rootRadius} identifies the exact valuation of $\zeta-1$. The final specialization theorem selects a bounded-order torsion point with controlled losses for a finite collection of polynomials. In \leanfile{Lib/PadicHodge/TateSen.lean}{the Tate--Sen file}, powers of the uniformizer $\zeta-1$ with exponent below $\varphi(p^h)$ have distinct valuation classes modulo $\Z$. The earlier graph report uses these classical local facts to evaluate individual conductor shells. We ask how to evaluate \emph{all conductors at once}. Our addition is a finite representation of the first surviving coefficient and a proof that its output gives the entire generating function. Neither RH nor an estimate for primes enters the graph argument.

\section{Valuations on a primitive conductor shell}
\label{sec:local-shell}

We prove Theorem~\ref{thm:main} by treating a finite Laurent polynomial before
specializing to graph Laplacians. Write
\[
 F(X)=p^g\sum_{i=1}^m c_i X^{a_i},\qquad
 a_i\in\Z^d,\quad c_i\in\Z\setminus\{0\},\quad
 \min_i v_p(c_i)=0 .
\]
Equal exponents have been collected. Fix $h\ge1$, $Q=p^{h-1}$,
$e=(p-1)Q$, and a primitive $p^h$-th root $\zeta$.
A point of joint conductor $p^h$ is $\zeta^u=(\zeta^{u_1},\ldots,
\zeta^{u_d})$, where at least one entry of $u\bmod p^h$ is a unit.
The simultaneous unit action on such tuples is free. Its orbits
have size $e$, and the valuation is constant on an orbit since
$\Q_p(\zeta)$ has a unique normalized $p$-adic valuation.
Thus a single representative contributes $e v_p(F(\zeta^u))$ to the
sum over the whole shell.

For each monomial choose $0\le n_i<pQ$ congruent to $a_i\cdot u$
modulo $pQ$, and write $n_i=r_i+Qb_i$ with $0\le r_i<Q$.
Set $A_{rb}=\sum_{i:(r_i,b_i)=(r,b)}c_i$. The unique representative
of the normalized evaluation $p^{-g}F(\zeta^u)$ in the power basis
$1,\zeta,\ldots,\zeta^{e-1}$ is $R(\zeta)$, where
\[
 R(Y)=\sum_{r=0}^{Q-1}\sum_{b=0}^{p-2}
 C_{rb}Y^{r+Qb},\qquad C_{rb}=A_{rb}-A_{r,p-1}.
\]
Indeed $\Phi_{p^h}(Y)=1+Y^Q+\cdots+Y^{(p-1)Q}$ eliminates the
last block of exponents. In particular $R=0$ detects exact vanishing.

\begin{lemma}[Integral orbit valuation]
\label{lem:orbit-val}
If $R\ne0$, let $\mu=\min_{C_{rb}\ne0} v_p(C_{rb})$.
Let $j$ be the first index whose coefficient in
$p^{-\mu}R(1+T)$ is nonzero modulo $p$. Then $0\le j<e$ and
\[
 e v_p(F(\zeta^u))=e(g+\mu)+j .
\]
\end{lemma}

\begin{proof}
Changing variables from $Y$ to $1+T$ preserves nonzero polynomials
over $\F_p$, so $j$ exists. The element $\pi=\zeta-1$ has
$v_p(\pi)=1/e$. In $p^{-\mu}R(1+\pi)$ the term of degree $j$
has valuation $j/e<1$. Every earlier nonzero coefficient is divisible
by $p$ and contributes valuation at least $1$, and each later term
has valuation strictly greater than $j/e$. The unique term of least
valuation gives the assertion, also when $p=2$ and $h=1$.
\end{proof}

The integer content is uniformly bounded: the sums defining the two terms in $C_{rb}$ use disjoint subsets of the coefficients, so $p^\mu\le\sum_i|c_i|$ when $R\ne0$. For a small example take $p=2,h=2$ and the square-grid symbol at $u=(1,1)$. The exponent polynomial $4-2X-2X^3$ reduces to $4$ modulo $X^2+1$. Thus $g=0,\mu=2,j=0$ and the orbit contributes $4$. Computing the content before collecting the terms would give the wrong answer.

\section{Computing all conductors with finitely many states}
\label{sec:lucas-machine}

Normalize each representative by making its first unit coordinate
equal to $1$ modulo $p^h$. In the chart with pivot $i_0$, lower-indexed
coordinates have units digit zero. The pivot's units digit is one, and
its other digits are zero. Every other digit is free. Reading the
coordinate vectors from the least significant end therefore
enumerates each orbit exactly once. At every stage there are
two operations on a new digit. Treating it as the top digit
emits a valuation at that conductor. Treating it as an ordinary digit
prepares the next stage.

After $k$ ordinary digits let $Q=p^k$, let $w\in\{0,\ldots,Q-1\}^d$
be the prefix, and decompose
\[
 a_i\cdot w=r_i+Q s_i,\qquad 0\le r_i<Q .
\]
Only the carries $s_i$ and a small amount of information about the $r_i$
are needed. For $0\le j<Q$ define a row vector over $\F_p$ by
\[
 B_j=\left(\binom{r_1}{j},\ldots,\binom{r_m}{j}\right)\pmod p .
\]
Process these rows in order of increasing $j$, retaining a row precisely
when it is outside the span of the earlier retained rows. Store the
ordered list of retained vectors and keep their integer indices in
separate registers. There are at most $m$ such registers.

\begin{lemma}[What the retained rows determine]
\label{lem:lucas-basis}
Two columns agree on the retained rows if and only if their residues
$r_i$ agree. For any linear functional on the rows, its first nonzero
value, if one exists, occurs on a retained row. The retained list after
one ordinary digit is computable from the previous list and the carries.
\end{lemma}

\begin{proof}
The full column associated to $r$ is the coefficient vector of
$(1+T)^r\in\F_p[T]$. Distinct $r<Q$ give different degrees. Agreement
on a basis implies agreement on all rows. A discarded row lies
in the span of earlier rows, so it cannot be the first detected by a
linear functional.

For a new coordinate digit vector $v\in\{0,\ldots,p-1\}^d$, perform
Euclidean division $s_i+a_i\cdot v=b_i+p s'_i$, $0\le b_i<p$.
The new residue is $r_i+Qb_i$ and the new carry is $s'_i$.
Over $\F_p$ one has $(1+T)^Q=1+T^Q$, hence
\[
 \binom{r_i+Qb_i}{j+Qt}
 \equiv \binom{r_i}{j}\binom{b_i}{t}\pmod p
 \qquad(0\le j<Q,\;0\le t<p).
\]
For each $t=0,\ldots,p-1$ in order, multiply every old retained row
coordinatewise by $(\binom{b_i}{t})_i$. Feed these candidates, in old
row order within the block, to the same greedy independence test.
The omitted old rows are combinations of earlier rows in their own
block after this multiplication. Inserting them into the scan would
therefore leave all decisions unchanged. A retained candidate coming
from an old index $j$ has new index $j+Qt$.
\end{proof}

\paragraph{Reading the top digit.}
Use the carries to compute $b_i$ as above. Partition the columns
according to their entries on the stored basis.
Lemma~\ref{lem:lucas-basis} identifies these groups with the occupied
values of $r$. Compute the integers $C_{rb}$ in each group, test for
vanishing, and then compute $\mu$ by Lemma~\ref{lem:orbit-val}.
Division by $p^\mu$ takes place \emph{after} the integer sums and
subtractions within each group. In top block $t=0,\ldots,p-2$ put
\[
 w_r(t)=\sum_{b=0}^{p-2}(C_{rb}/p^\mu)\binom bt\pmod p .
\]
At index $j_0+Qt$, the coefficient to be tested modulo $p$ is
$\sum_r w_r(t)\binom{r}{j_0}$. This is a linear functional on
the old rows, using one coordinate per group. Testing the stored rows
in order finds its earliest nonzero value. Stop at the first successful
block. If the successful old index is $j_0$, the orbit contribution is
\[
 \big((p-1)(g+\mu)+t\big)Q+j_0 .
\]
All choices except the numerical values of $Q,j_0$ depend on the finite
state and this last digit.

\paragraph{A reduction of the state space.}
After processing $k$ ordinary digits, test the normalized coefficient
functional $B_j\mapsto\sum_i c_i\binom{r_i}{j}$. Suppose its first
nonzero value is at $j_*<p^k$. For every subsequent conductor the
ordinary exponent polynomial at $1+T$, reduced modulo $(p,T^{p^k})$,
equals $\sum_i c_i(1+T)^{r_i}$ in this quotient.
Indeed further exponent digits multiply these terms by powers of
$1+T^{p^k}$. Final cyclotomic reduction preserves the same truncation,
because $\Phi_{p^h}(1+T)\equiv T^{\varphi(p^h)}\bmod p$ and
$\varphi(p^h)\ge p^k$ for $h>k$. Consequently every continuation has
$\mu=0$ and $j=j_*$, including arbitrary choices for its final digit.
The implementation stores only $Q$ and $j_*$ in a new state with the
same chart. Future ordinary digits replace $Q$ by $pQ$ and leave
$j_*$ fixed.

\paragraph{Termination and weights.}
The control state consists of the pivot, a first-digit flag, the
carries, and the ordered list of row vectors. The register values are
stored separately. If $H_i=\sum_\ell |a_{i\ell}|$, then
$|s_i|=|\lfloor a_i\cdot w/Q\rfloor|\le H_i$ for every prefix.
An explicit upper bound on the number of control states, including the
reduced states, is
\[
 N=d+2d\left(\prod_{i=1}^m(2H_i+1)\right)\sum_{s=0}^m p^{ms}.
\]
A breadth-first search enumerates the reachable states and tests every
allowed final digit. It therefore also decides whether a nontrivial
$p$-power torsion zero exists. One can reject such inputs or, for a
sum taken only over nonzero evaluations, give a zero evaluation
weight zero while continuing the search. Rejection is appropriate
for an ordinary product.

Each ordinary state has registers $(Q,j_1,\ldots,j_s)$ with $s\le m$.
Each ordinary digit induces an integer linear map between registers by
$Q'=pQ$, $j'_a=j_b+tQ$. Summing the maps for digits with the same target
gives the corresponding block of a matrix $M$. In each state the sum
of the output forms over its legal top digits gives the corresponding
block of a row $f$. Initially each chart has zero carries, the single
row $(1,\ldots,1)$, and registers $(1,0)$. Let $u$ be the concatenated
initial vector, zero outside these initial states. Matrix iteration
adds the registers of all prefixes at each state. By induction on
prefix length and the unique chart representation, the shell sum at
conductor $p^h$ is exactly
\[
 b_h=fM^{h-1}u .
\]
The matrix dimension $L$ is the sum of the register counts of the
reachable states. In particular $L\le(m+1)N$. The implementation
uses the dimension obtained from its search.

\section{Spanning trees and the first level of the eventual law}
\label{sec:exact-gf}

Let $G$ be a finite connected undirected multigraph with integral
voltages $a_e\in\Z^d$, one orientation per edge. Its Laurent Laplacian
adds $1$ to both endpoint diagonals and subtracts $X^{a_e},X^{-a_e}$
in the two corresponding off-diagonal positions, interpreting all four
entries at the same position for a loop. Take $F$ to be its determinant.
The finite lifts have deck group $(\Z/p^n\Z)^d$ and we assume they
are connected. Character decomposition of the lifted Laplacian
gives one block per character. The trivial block is the base Laplacian.
The matrix-tree theorem applied to the base and lift cancels their
common vertex-count factor and leaves the deck-group order in the
denominator. Hence, writing $v_n=v_p(\kappa_n)$ and
$v_0=v_p(\kappa(G))$,
\[
 v_n=v_0-dn+\sum_{h=1}^n b_h .
\]
The base cofactor and Laurent determinant are explicit finite input
calculations. The same argument works for graphs with leaves and for
cycles. A connected base has a connected lift at a given level exactly
when every nontrivial character block there is invertible, because
the multiplicity of the zero eigenvalue counts components.

With $z$ a formal variable, summing the last display yields
\begin{equation}
 \sum_{n\ge0}v_n z^n=
 \frac{v_0}{1-z}-\frac{dz}{(1-z)^2}
 +\frac{z}{1-z}\,f(I-zM)^{-1}u .
 \label{eq:machine-resolvent}
\end{equation}
This identity computes the entire sequence. If
$\chi_M(x)=x^L+c_1x^{L-1}+\cdots+c_L$, then
$\det(I-zM)=1+c_1z+\cdots+c_Lz^L$. The adjugate shows that the
numerator of $f(I-zM)^{-1}u$ over this determinant has degree at most
$L-1$. Multiplying that denominator by the first $L$ shell coefficients
of the series and keeping degrees below $L$ gives its exact numerator,
including when $M$ is singular. This is the characteristic-polynomial
procedure used in the accompanying implementation.

For the Greenberg polynomial put
\[
 D(z)=(1-p^d z)\prod_{i=0}^{d-1}(1-p^i z)^2 .
\]
An eventual law $v_n=P(p^n,n)$ with $\deg P\le d$ and
$\deg_Y P(X,Y)\le1$ implies that the reduced denominator of
\eqref{eq:machine-resolvent} divides $D(z)$. Conversely, whenever it
divides $D(z)$, polynomial division and partial fractions give
\[
 \sum_{n\ge0}v_nz^n
 =H(z)+\sum_{i=0}^d\left(\frac{A_i}{1-p^i z}
            +\frac{B_i}{(1-p^i z)^2}\right),\qquad B_d=0 .
\]
Thus $P(X,Y)=\sum_i[A_i+(Y+1)B_i]X^i$ and
$v_n-P(p^n,n)=[z^n]H$. If $H=0$ the least onset is zero.
Otherwise it is $1+\deg H$. Uniqueness follows since the largest
exponential base in a nonzero sum $\sum_i(a_i+b_i n)p^{in}$ dominates
its lower bases on a tail. DuBose--Vallières guarantee the denominator
shape for base graphs with no leaves and nonzero Euler characteristic.
To pass to an arbitrary connected base admitting such a connected tower,
repeatedly remove pendant edges. Their lifts are pendant edges and belong
to every spanning tree. If the remaining core is a cycle, its fundamental
group is cyclic, so a connected tower with deck groups as above has
$d=1$. A cycle of length $b$ then lifts to a cycle of length $bp^n$,
with $v_n=v_p(b)+n$ (including a single loop or a double edge).
A tree admits no connected nontrivial finite cover. These cases exhaust
the missing base graphs. There is also a numerical upper bound on the required conductor precision. Formula~\eqref{eq:machine-resolvent} has numerator and denominator degrees at most $L+2$. Consequently $\deg H\le L+2$ if $H\ne0$, and the least onset is at most $L+3$. The characteristic-polynomial computation itself requires only the first $L$ shell sums after the reachable matrix has been built. Computed divisibility
also certifies the shape directly in each example. For the general
Laurent sum, the characteristic-polynomial procedure produces its
exact rational series independently of that existence theorem.

The integer-voltage formulation also covers rational voltages in
$\Q\cap\Z_p$. Multiplying every voltage by a common denominator prime
to $p$ induces compatible deck-group automorphisms at every level,
so the finite graphs are isomorphic to those with the cleared
integral voltages.

\section{Five complete computations}
The five examples in DuBose--Vallières \S7 provide a check against published computations. The prior ProveIt packet proves the five formulas by separate arguments. Supplying only the prime and voltages to the same program in all five cases gives the following output. Write $X=p^n$. The two right-hand columns give the least start including level zero and the total dimension of the integer linear representation.
\begin{center}\small
\begin{tabular}{@{}ccp{.31\linewidth}lrr@{}}
\toprule
Case & $p$ & Steps & $P(X,n)$ & $N$ & Dim.\\
\midrule
1&2&$e_1,e_2$&$(2n+4)X-6n-1$ &1&25\\
2&2&$e_1,e_1,e_2,e_2$&$X^2+(2n+4)X-6n-2$ &1&25\\
3&2&$(1,5),(0,3),(1,2),(0,1)$&$(n+33/4)X-4n-1$ &4&32\\
4&3&$e_1,e_2$&$4X-2n-4$ &0&8\\
5&3&$(1,0),(2,3),(1,1)$&$(20/3)X-2n-8$ &1&16\\
\bottomrule
\end{tabular}
\end{center}
In case 3, the four initial values including level zero are $0,5,19,65$. The finite correction is
\[
 E(z)=-29/4-(17/2)z-13z^2-12z^3.
\]
For example, the stable polynomial predicts $77$ at level three. The exact value is $65$. The degree-three correction proves that level four is the first stable level; its derivation uses all future conductors through the transfer matrix.

The file \texttt{dv\_examples.json} records each input, rational generating function, tail polynomial, finite correction, and the values through level ten.
The program \texttt{check\_resultants.py} reruns all five examples and checks small Laurent polynomials at $p=2,3,5$ against exact cyclotomic resultants. Other checks in the archive use integer Laplacian cofactors and determinants of Grover transition matrices. The proofs here are ordinary mathematical proofs. The programs supply reproducible exact calculations, and the new theorems have not been formalized in Lean.

\subsection*{A two-vertex example for every prime}
Consider the three-edge graph in \texttt{honeycomb.json}. Its two vertices are joined by edges of voltages $0,e_1,e_2$. Its covers are the periodic honeycomb graphs. The base has three spanning trees. The determinant of its two-by-two twisted Laplacian is
\[
 F(X,Y)=9-(1+X+Y)(1+X^{-1}+Y^{-1})
        =H(X)+H(Y)+H(X/Y),\qquad H(X)=2-X-X^{-1}.
\]
At $p=2,3$ the program yields, respectively, the following all-level formulas, with zero finite correction in both cases:
\[
 T_n=(3n+9)2^n-8n-9,\qquad T_n=5\cdot3^n-2n-4 .
\]
The register dimensions are 44 and 12. In particular the ternary constant term takes account of the base tree count.

For $p\ge5$ a short calculation gives the full family. At every level,
\[
 T_n=A_p\,\frac{p^n-1}{p-1}-2n,\qquad
 A_p=\begin{cases}
 2(p+5),&p\equiv1\pmod3,\\
 2(p+1),&p\equiv2\pmod3 .
 \end{cases}
\]
Here is a direct proof. In a primitive direction $(a,b)$ modulo $p$, the first even moment is $a^2+b^2+(a-b)^2=2(a^2-ab+b^2)$. If it is nonzero, the first unit degree is two. Otherwise $ab\ne0$ and $b/a$ is a primitive sixth root. The three squared projections $a^2,b^2,(a-b)^2$ are $a^2$ times the three cube roots of unity. Their first two power sums vanish and their third is $3a^6\ne0$. There are two such projective directions for $p\equiv1\bmod3$ and none for $p\equiv2\bmod3$.

To identify the first unit degree in the exceptional directions, use $X=\exp(az),Y=\exp(bz)$ formally through degree six modulo $p$; here $p\ge7$, so all factorials in this expansion are units. The coefficients at even degree $2j$ are $-2/(2j)!$ times the $j$th power sum of the squared projections. The change from $z$ to $\exp(z)-1$ preserves the order, so the first unit degree is six. In a cyclotomic specialization, the integral Laurent expansion in $T=\zeta-1$ has its first two coefficients identically zero. Thus even at $p=7,h=1$, when degree six equals the ramification index, every lower nonzero term divisible by $p$ has normalized order at least $e+2>6$. These unit degrees are exactly the orbit orders at every conductor. Each projective direction has $p^{h-1}$ normalized lifts at level $h$. Summing them gives $A_p$ per direction set at level one and the formula above by geometric summation. This proof specializes the moment method already present in the ProveIt packet~\cite{incoming}.

\section{The first stable level can be prescribed}

Fix a prime $p\equiv3\pmod4$ and an integer $k\ge1$. On
$(\mathbb Z/p^n\mathbb Z)^2$ take one step in each of the directions
$p^k e_1$ and $e_2$, and $p$ copies of the step $e_1$. Each step from a
vertex contributes one undirected edge. Parallel edges count separately
and loops contribute zero to the Laplacian. These connected towers
all have the same number $p+2$ of base loops for fixed $p$.
Let $T_n$ be the $p$-adic valuation of the spanning-tree count.

\begin{theorem}
Put $s=\min(k,n)$ and $C=2(p+1)/(p-1)$. At every level $n\ge1$,
\[
 T_n=2s p^n+C(p^n-p^s)+(p^s-1)+2s-2n .
\]
The eventual polynomial is
\[
 P_k(X,Y)=(2k+C)X-2Y+(1-C)p^k+2k-1.
\]
It gives $T_n=P_k(p^n,n)$ from level $k$ onward. For $k\ge2$ it fails
at level $k-1$, so $k$ is the exact first stable level.
\end{theorem}

\begin{proof}
Write $A(U)=2-U-U^{-1}=-(U-1)^2/U$.
For a pair $(U,V)$ of $p$-power roots of unity of joint conductor $p^j>1$,
let $e_j=(p-1)p^{j-1}$. In a common primitive coordinate $\zeta$ of order
$p^j$, write $U=\zeta^a$, $V=\zeta^b$, with $a,b$ not both divisible
by $p$. The reduction of
$(A(U)+A(V))/(\zeta-1)^2$ modulo $\zeta-1$ is $-(a^2+b^2)$. Since
$-1$ is a nonsquare modulo $p$, this is nonzero. Consequently
\[
 v_p(A(U)+A(V))=2/e_j\le1 .
\]
The eigenvalue at a character $(X,Y)$ is
$A(X^{p^k})+A(Y)+pA(X)$. The final term has valuation strictly above one,
with the convention $v_p(0)=+\infty$. Thus whenever
$(X^{p^k},Y)\ne(1,1)$, the joint conductor of this image alone determines
the valuation.

The map $(X,Y)\mapsto(X^{p^k},Y)$ has fibers of size $p^s$ and image
isomorphic to $C_{p^{n-s}}\times C_{p^n}$, where $C_m$ denotes a cyclic
group of order $m$. For $1\le j\le n-s$ its exact conductor-$p^j$ shell
has $p^{2j}-p^{2j-2}$ elements. Their aggregate valuation is therefore
$2(p+1)p^{j-1}$. For $n-s<j\le n$ there are $p^{n-s}e_j$ elements in
the shell, contributing $2p^{n-s}$. Summing these shell contributions
and multiplying by $p^s$ gives $C(p^n-p^s)+2s p^n$.
On the kernel with the trivial character removed, the eigenvalue is
$pA(X)$ for each nontrivial $p^s$-th root $X$. There are $p^s-1$ such
roots, and
$\prod_{X\ne1}(1-X)=p^s$ gives $\sum_{X\ne1}v_p(A(X))=2s$.
The matrix-tree denominator contributes $-2n$. This proves the formula.

For $n\le k$ it reads $T_n=(2n+1)p^n-1$, and for $n\ge k$ it equals
$P_k(p^n,n)$. In particular the discrepancy one level before $k$ is
\[
 T_{k-1}-P_k(p^{k-1},k-1)=(p+1)p^{k-1}-2>0 .
\]
Distinct polynomials in $p^n,n$ cannot agree on an infinite tail:
the term of highest degree in $p^n$ dominates all lower exponential
orders, even with polynomial factors in $n$. This establishes the
claimed onset for the unique eventual polynomial.
\end{proof}

The growing onset is visible directly in the formal specializations.
For primitive $(a,b)\in\mathbb Z_p^2$, substitute $X=(1+z)^a$ and
$Y=(1+z)^b$ in the Laplacian polynomial and reduce modulo $p$.
The first nonzero degree is
\[
 2p^{\min(k+v_p(a),v_p(b))}.
\]
Indeed the two summands that survive reduction have these respective
initial orders, and in an equal-order tie the same anisotropic sum of
squares prevents cancellation. The maximum is exactly $2p^k$, attained
at $(1,0)$. Thus even within this family with no identically zero
primitive specialization, neither the maximal degree required to see a
unit nor the onset admits a bound in terms of $p$, dimension and base
edge count alone. Replacing $p^k e_1$ by the zero step gives the law
$(2n+1)p^n-1$ at every level, despite agreement of the voltage data
modulo $p^k$.

For $p=3$ and $k=3$ the stable law is $10\cdot3^n-2n-76$ and the first
five values are $8,44,188,726,2344$. The script
\texttt{verify\_delayed\_tower.py} checks the eigenvalues by exact
cyclotomic reduction. It independently constructs integer Laplacian
cofactors for its examples with at most 81 vertices.

\paragraph{A version with distinct steps.}
For $p\ge7$ in the same congruence class, replace the $p$ copies of
$e_1$ by $(1+i p^k)e_1$ for $0\le i<p$. Every valuation stays the same.
Indeed, writing $Z=X^{p^k}$, the replacement contributes
$2p-X\Phi_p(Z)-X^{-1}\Phi_p(Z^{-1})$. At $Z=1$ this is $pA(X)$.
If $Z$ has order $p$ the two cyclotomic factors vanish. If its order is
larger, each has valuation $(p-1)v_p(Z-1)$ by
$\Phi_p(Z)=(Z^p-1)/(Z-1)$. Hence for $Z\ne1$ the new contribution has
valuation strictly greater than $2v_p(Z-1)$, which in turn is at least
$v_p(A(Z)+A(Y))$. For $n>k$ all steps are nonzero and distinct up to
sign, so the graphs at those levels are simple.

\section{Even the leading coefficient can change after an arbitrary level}
\label{sec:seven-loop}
Fix $p=3$. Take three copies of each of $(1,0)$ and $(0,1)$ and one
additional step $(t,0)$. The number of steps stays at seven as $t$ varies.
Write $T_n(t)$ for the valuation of the spanning-tree count and
$k\ge1$. At $t=0$ the last step gives loops. Put
\[
 B(n)=9^n+4\cdot3^n-2n-5,\qquad
 D_k(n)=(2(n-k)+1)3^{n+k}+(2(n-k)+4)3^k-2n-5.
\]
Then $T_n(0)=B(n)$ for every positive $n$, and
\[
 T_n(3^k)=
 \begin{cases}
 B(n),&1\le n\le k,\\
 D_k(n),&n\ge k.
 \end{cases}
\]
Both expressions agree at $k$. The least positive starting level for the
eventual formula $D_k$ is exactly $k$.

\paragraph{Proof.}
Index a character by its two coordinate roots of unity. Let $a,b$ be their
conductor exponents, with exponent $0$ for the trivial root. For a
nontrivial character write $h=\max(a,b)$ and express the roots as
$z^u,z^v$ for $z$ primitive of order $3^h$. At least one of $u,v$ is
a unit mod $3$. With $H(w)=2-w-w^{-1}$ and $\pi=z-1$, the grid symbol
$F=H(z^u)+H(z^v)$ satisfies
$F/\pi^2\equiv-(u^2+v^2)\pmod\pi$. The nonzero squares modulo $3$
are all $1$, so this residue is nonzero. Consequently $v_3(F)=
2/\varphi(3^h)=3^{1-h}$.
If $a\le k$, the full symbol equals $3F$ and has valuation
$1+3^{1-h}$. If $a>k$, the extra summand $H(z^{3^k u})$ has valuation
$3^{1-a+k}\le1$, strictly below that of $3F$, and gives the valuation
of the sum.

For $n\le k$ there are $9^n-1$ nontrivial characters, and the
$8\cdot9^{h-1}$ characters at conductor $h$ contribute
$8\cdot3^{h-1}$ from the grid valuation. Summing and subtracting the
matrix-tree correction $2n$ gives $B(n)$. This argument holds for every
$n$ at $t=0$. For $n\ge k$, the rectangle $a\le k$ contains
$3^{n+k}-1$ nontrivial characters. Their unit contribution is this
count. Their grid valuations sum to $4(3^k-1)$ for $h\le k$ and to
$2(n-k)3^k$ for $h>k$, since then $b=h$. Outside the rectangle,
the $2\cdot3^{a-1+n}$ characters with given $a>k$ contribute
$2\cdot3^{k+n}$. Adding these sums and subtracting $2n$ gives
$D_k(n)$. For $k\ge2$ the difference at the preceding level is
\[
 B(k-1)-D_k(k-1)=4\cdot3^{2k-2}-2\cdot3^{k-1}>0.
\]
This proves minimality, including $k=1$ by the positive-level convention.

\paragraph{Leading coefficients and finite input precision.}
Write $X=3^n$. Denote the coefficient of $X^2$ by $\mu$ and that of
$nX$ by $\lambda$ in the eventual polynomial. At $t=0$ the pair is
$(1,0)$. At $t=3^k$ it is $(0,2\cdot3^k)$. Thus the degree in $X$
drops from two to one, and $\lambda$ is unbounded along a sequence for
which all voltage digits stabilize. In particular even $\mu$ is
discontinuous at $0$ on the space of $3$-adic inputs. Any digit-query
algorithm halting at $0$ has read only finitely many digits. A
sufficiently large $3^k$ yields the same answers to all these queries,
yet the required leading coefficient differs. This rules out a total
exact digit-oracle algorithm even for this single coefficient. The same obstruction holds for
finite programs computing the digits, even with a promise that their
values are ordinary integers. Given a machine, output the residues of
$3^j$ if it first halts at time $j\ge1$, and the residues of zero if it
never halts. Computing the residue modulo $3^r$ needs at most $r$
simulation steps. An algorithm for $\mu$ on such programs would decide
halting. A finite integer encoding distinguishes these inputs.

The independent cyclotomic check in \texttt{verify\_leading.py} computes
both sides across the change of formula. For $k=3$, the successive
valuations at levels $1$ through $5$ are $14,108,826,6710,33006$.

\paragraph{A fixed simple base graph.}
The same formulas hold with a simple graph at every level. Replace each
of the seven base loops by a cycle of length four, using separate internal
vertices for each cycle. Put the original voltage on one oriented edge
of that cycle and zero on the other three. This replaces every edge in
each lifted graph by a path of length four. For a connected multigraph
with $E$ edges and $V$ vertices, the resulting spanning-tree count is
$4^{E-V+1}$ times the original count. Indeed, the intact paths project
to a spanning tree, and on each remaining path one of the four edges
must be removed. The argument also applies to subdivided loops.
The factor is a $3$-adic unit, so the least onset $k$, the discontinuity
of $\mu$, and the unbounded values of $\lambda$ all persist on covers
of this single simple base.

\section{Exact spectral tests for Grover walks}
A Grover walk on a finite graph acts on its directed edges. From an arc ending
at $v$, the amplitude to each outgoing arc is $2/\deg(v)$, with $1$
subtracted for the reversed arc. A loop has two directed arcs, exchanged by
reversal. On an abelian cover, Fourier transformation along the deck group
splits the transition matrix into blocks $U(\xi)$, one for each character
$\xi$. For integer voltages the block entries are rational multiples of
Laurent monomials in the coordinates of $\xi$: a deck translation contributes
its character value, and an opposite voltage contributes the inverse.

\subsection*{Deciding spectral occurrence and computing determinants}
Given $P\in\mathbb Q[A]$, form the rational Laurent polynomial
$F_P(X)=\det P(U(X))$ on the finite base arc space.
At level $n$ one has the exact factorization
\[
                \det P(U_n)=\prod_{\xi\in\mu_{p^n}^d}F_P(\xi).
\]
This is the direct arc-space version of the spectral resultant factorization
of Akahori, Hayashi and Suzuki \cite[Theorems 3.3 and 4.3]{AHS}.
Using the full block retains the eigenvalues $1$ and $-1$ as well.
If $F_P$ is identically zero or $F_P(1)=0$, a base-level eigenvalue
already answers the occurrence question. Otherwise choose a positive
integer $b$ with $G=bF_P$ integral. A reachable exact-zero test for $G$
decides whether a root of $P$ occurs as an eigenvalue at any finite layer.
If it does, a shortest accepted digit word gives the first such level.
If it does not, the valuation sequence is the sum of the nontrivial
character valuations for $G$, plus $v_p(G(1))-p^{dn}v_p(b)$.
Thus the finite automaton gives its exact generating function. Monsky's
torsion-sum theorem \cite{monsky} supplies the eventual polynomial shape
in the nonvanishing case; clearing $b$ adds the displayed degree-$d$
term. The computation recovers that polynomial and its least starting
level. The leading spectral asymptotics in \cite[Theorem 4.3]{AHS}
allow more general coefficient and voltage data.

\subsection*{An all-level formula on the ternary square grid}
Take the square torus with side $3^n$, $N=9^n$ vertices, and the four
directional arcs at each vertex, including at $n=0$. Set
$S(X,Y)=4-X-X^{-1}-Y-Y^{-1}$. The characteristic polynomial is
\[
 \det(AI-U_n)=(A^2-1)^N
  \prod_{X,Y\in\mu_{3^n}}\left((A-1)^2+\tfrac A2 S(X,Y)\right).
\]
Indeed each of the $N$ character blocks has the two factors $A^2-1$ and
$A^2+1-(A/2)(X+X^{-1}+Y+Y^{-1})$ by the Grover spectral
factorization, giving the displayed identity.

\begin{theorem}
For every $n\ge0$ and integer $a\ne\pm1$,
\[
 v_3(\det(aI-U_n))=
 \begin{cases}
 0,& a\equiv0\pmod3,\\
 m9^n+4\cdot3^n+2m-4,& a\equiv1\pmod3,\quad m=v_3(a-1),\\
 v_3(a+1)9^n,&a\equiv-1\pmod3 .
 \end{cases}
\]
At $A=1$ the zero has order $N+2$, with first nonzero Taylor coefficient
of valuation $4(3^n-1)$. At $A=-1$ the zero has order $N$ and the
first nonzero Taylor coefficient is a $3$-adic unit.
\end{theorem}
\begin{proof}
If $(X,Y)$ has joint conductor $3^h>1$, the ternary sum of squares argument
above gives $v_3(S)=3^{1-h}$. There are $8\cdot9^{h-1}$ such pairs,
so the sum of $v_3(S)$ over nontrivial pairs is $4(3^n-1)$. At $(1,1)$
the value of $S$ is zero.

If $a\equiv1\pmod3$, the term $(a/2)S$ in every nontrivial block has
valuation at most one, and $(a-1)^2$ has valuation $2m\ge2$.
The trivial block contributes $2m$ and the factor $(a^2-1)^N$ contributes
$mN$. For the two other residue classes $(a-1)^2$ is a unit while $aS/2$
is in the maximal ideal in each character field. Only $(a^2-1)^N$
can then contribute. At $A=1$ the trivial quadratic block contributes
a double zero; the leading coefficient is $2\prod_{(X,Y)\ne(1,1)} S(X,Y)$.
At $A=-1$ all quadratic blocks are units.
\end{proof}

For example, at $a=4$ the successive values for $n=0,1,2$ are
$3,19,115$. The script \texttt{verify\_grover.py} checks exact determinants
of the full arc matrices for $n=0,1$ and of the vertex factor through $n=2$.

\subsection*{Periodic states on an odd rectangular grid}
Remark 6.5 of \cite{AHS} proposes using a reduced polynomial derived from
$A^r-1$ to study periodic eigenvalues after removing the factors at $\pm1$.
Here the entire periodic subspace can be counted. Let the side lengths
$L,M$ be odd, $N=LM$ and $g=\gcd(L,M)$. In addition to the eigenvalues
$1$ and $-1$, of multiplicities $N+2$ and $N$, respectively, the
root-of-unity eigenvalues are precisely the nontrivial $g$-th roots,
each of multiplicity four. Thus the periodic subspace has dimension
$2LM+4g-2$, and the restriction of the walk to it has exact period $2g$.

To verify the only delicate point, a root of unity $w$ in a quadratic
block must satisfy
\[
              2(w+w^{-1})=X+X^{-1}+Y+Y^{-1}.
\]
This is a special case of the rational-angle cosine equation treated by
Speyer \cite{speyer}. On an odd rectangular grid a short residue argument
suffices. Reduce in the common cyclotomic field at a prime above $2$.
The residue of $Y$ must solve the quadratic with roots the residues of
$X,X^{-1}$. Distinct roots of odd order stay distinct under this reduction:
if a nontrivial odd-order root $u$ reduced to $1$, its vanishing geometric
sum would reduce to its odd order, which is nonzero. Hence $Y=X$ or
$Y=X^{-1}$, and substitution forces $w=X$ or $X^{-1}$. For a nontrivial
$g$-th root $w$, the four pairs in $\{w,w^{-1}\}^2$ each contribute one
simple root. For $w=1$ only $(1,1)$ contributes the extra double root,
and for $w=-1$ no quadratic block contributes. The unitary walk is
diagonalizable, so these algebraic multiplicities count dimensions.

\section{Further research questions}
\begin{enumerate}
\item For fixed $p,d$ and fixed number of voltage vectors, what is the largest possible onset for vectors of bit length at most $b$? The five-loop examples give growth at least linear in $b$ at $p=3$. The finite-state bound in the proof is much larger. Even the correct order for two-dimensional bouquets is worth determining.
\item Can one compute the smallest rational linear representation directly from the exponent vectors, without constructing all reachable digit states? The five DuBose--Vallières examples have far fewer poles than registers. A structural description of the canceled factors could yield a practical improvement for Laplacians with many terms.
\item Which exact input representations beyond rational voltages support a terminating algorithm? Simultaneously replacing all vectors by their images under a specified deck-group automorphism preserves every graph. More general algebraic $p$-adic inputs would need additional control of carries or an alternative way of recognizing coincident exponents. The digit-oracle obstruction explains why an input model matters here.
\end{enumerate}

\begin{thebibliography}{9}\small
\bibitem{repo} L. Sernau and collaborators, \emph{GRH for Dirichlet $L$-functions, in Lean 4}, based on work of L. Alpöge with Claude.
\href{https://github.com/lsernau/grh-lean/tree/72fead102f663d5c7a407964a13dc0cf9d977207/grh}{Readable edition and verification notes};
\leanfile{Statements.lean}{statements};
\leanfile{Main.lean}{five public theorems}. Lean 4.35.0, Mathlib \texttt{56a00585}.
\bibitem{incoming} \emph{Cyclotomic valuations in graph towers and progression-free sets over finite rings}, prior ProveIt packet
\href{https://github.com/VladimirReshetnikov/ProveIt/blob/1667e16cee9d4a68961e201e7288b8ecdac895c6/docs/incoming/grh_method_transfers_graph_towers_and_finite_rings.zip}{graph towers and finite rings}. See the graph chapter and \texttt{endmatter.tex}, first research question. Git blob \texttt{8a901afb0086eb6574bb0d07a9081b3fba3d9907}.
\bibitem{DV} S. DuBose and D. Vallières, \emph{On $\Z_\ell^d$-towers of graphs}, Algebraic Combinatorics \textbf{6} (2023), 1331--1346. \href{https://doi.org/10.5802/alco.304}{Theorem 6.2 and Section 7}.
\bibitem{monsky} P. Monsky, \emph{On $p$-adic power series}, Math. Ann. \textbf{255} (1981), 217--227, Theorem 5.6. See also S. Tateno and J. Ueki, \emph{The Iwasawa invariants of $\Z_p^d$-covers of links}, J. Lond. Math. Soc. \textbf{111} (2025), e70183, \href{https://doi.org/10.1112/jlms.70183}{Proposition 3.9}.
\bibitem{AHS} J. Akahori, T. Hayashi and R. Suzuki, \emph{Iwasawa-type spectral resultant growth laws for Grover walks on graph towers}, Ann. Math. Québec (2026). \href{https://doi.org/10.1007/s40316-026-00290-1}{Theorem 3.3, Corollary 4.6, Theorem 6.2 and Remark 6.5}.
\bibitem{speyer} D. E. Speyer, \emph{Answer on rational-angle solutions to a cosine product equation}, MathOverflow (2023). \href{https://mathoverflow.net/questions/458827}{Dyadic argument and the equivalent sum equation}.
\end{thebibliography}
\end{document}
